% Provenance (never printed): see sources/outline_sources.md and sources/book1/BRIEFS.md (g8/04).
\chapter{Searching and Sorting}\label{ch:g8:searching-and-sorting}

To find the word \emph{otter} in a paper \omterm{def:g8:dictionaries:dictionary}{dictionary}, nobody starts at the
first page and reads every word. You open it near the middle, see that you
are in the \emph{m}s, and go on in the second half; a few openings later the
word is in front of you. That works because the words are in alphabetical
order. This chapter teaches a \omterm{def:g1:what-is-a-computer:computer}{computer} both ways of searching, and two ways
of putting a \omterm{def:g7:lists:list}{list} in order first.

\begin{recall}
A \omterm{def:g7:lists:list}{list} keeps \omterm{def:g7:lists:list}{items} in order, each at a position, its \omterm{def:g7:lists:list}{index}, from 0
(\cref{ch:g7:lists}); a \omterm{def:g7:functions:function}{function} takes \omterm{def:g7:functions:parameter}{parameters} and returns a \omterm{def:g6:variables-and-input:value}{value}
(\cref{ch:g7:functions}). The guessing game of \cref{ch:g6:while-loops} asks
``higher or lower?'' until the guess is right. $n$ \omterm{def:g5:bits-and-bytes:bit}{bits} make $2^n$ patterns,
and each extra \omterm{def:g5:bits-and-bytes:bit}{bit} doubles their number (\cref{ch:g8:binary-and-hexadecimal}).
\end{recall}

\section{Searching item by item}

\begin{definition}[Linear search]\label{def:g8:searching-and-sorting:linear-search}
A \emph{linear search}\index{linear search} looks at the \omterm{def:g7:lists:list}{items} of a \omterm{def:g7:lists:list}{list} one
after the other, from the first, until it finds the one it looks for or
reaches the end.
\end{definition}

The \omterm{def:g7:functions:function}{function} returns the position of the target, or $-1$ when it is not in
the \omterm{def:g7:lists:list}{list}, a position that cannot exist. It also counts its \omterm{def:g6:if-else:comparison}{comparisons}, to
measure its work.

\omcode{code/grade-8/04-searching-and-sorting/listings/linear_search.py}{1}{14}{\omterm{def:g8:searching-and-sorting:linear-search}{Linear
search}, counting its \omterm{def:g6:if-else:comparison}{comparisons}.}\label{lst:g8:searching-and-sorting:linear}

\Needspace*{7\baselineskip}
\omoutput{code/grade-8/04-searching-and-sorting/outputs/linear_search.txt}{Position
and \omterm{def:g6:if-else:comparison}{comparisons}.}\label{out:g8:searching-and-sorting:linear}

The \omterm{def:g7:functions:function}{function} returns two \omterm{def:g6:variables-and-input:value}{values}, written with a comma; Python prints them in
brackets. Finding 12 takes one \omterm{def:g6:if-else:comparison}{comparison}, because it is first; 13 takes 7.
Looking for 7, which is not there, takes 8: every \omterm{def:g7:lists:list}{item} is looked at. This is
the \emph{worst case}: for a \omterm{def:g7:lists:list}{list} of a million \omterm{def:g7:lists:list}{items}, a million \omterm{def:g6:if-else:comparison}{comparisons}.

\begin{omfigure}
\begin{tikzpicture}[x=0.9cm, font=\small]
  \foreach \v [count=\i from 0] in {12, 9, 15, 11, 9, 14, 13, 10} {
    \draw[black!60, fill=ocBlue!8] (\i,0) rectangle ++(0.9,0.8);
    \node at (\i+0.45,0.4) {\v};
    \node[font=\scriptsize, black!60] at (\i+0.45,-0.25) {\i};
  }
  \draw[ocMint!80!black, very thick, fill=ocMint!25] (6,0) rectangle ++(0.9,0.8);
  \node at (6.45,0.4) {13};
  \foreach \i in {0,...,5} {
    \draw[-{Stealth[length=1.5mm]}, ocRed] (\i+0.5,1.0) to[bend left=50] (\i+1.4,1.0);
  }
  \node[anchor=west, font=\footnotesize] at (8.2,0.4) {13 found after 7 comparisons};
\end{tikzpicture}

{\small \omterm{def:g8:searching-and-sorting:linear-search}{Linear search} for 13: one cell after the other, from the left.}
\end{omfigure}

\section{Searching a sorted list by halving}

\begin{definition}[Sorted list]\label{def:g8:searching-and-sorting:sorted}
A \omterm{def:g7:lists:list}{list} is \emph{sorted}\index{sorted list} (in increasing order) when each
\omterm{def:g7:lists:list}{item} is smaller than or equal to the next one.
\end{definition}

\begin{definition}[Binary search]\label{def:g8:searching-and-sorting:binary-search}
A \emph{binary search}\index{binary search} looks for an \omterm{def:g7:lists:list}{item} in a \omterm{def:g8:searching-and-sorting:sorted}{sorted
list} by comparing it with the middle \omterm{def:g7:lists:list}{item}, then going on in the half where it
must be, and halving again until it is found or no \omterm{def:g7:lists:list}{item} is left.
\end{definition}

This is the good strategy of the guessing game: guess 50 between 1 and 100;
``higher'' leaves 51 to 100, so guess 75; and so on. The \omterm{def:g4:programs-and-operating-system:program}{program} keeps two
positions, \code{low} and \code{high}: the target, if it is there, lies
between them.

\omcode{code/grade-8/04-searching-and-sorting/listings/binary_search.py}{1}{21}{\omterm{def:g8:searching-and-sorting:binary-search}{Binary
search} between \code{low} and \code{high}.}\label{lst:g8:searching-and-sorting:binary}

\Needspace*{8\baselineskip}
\omoutput{code/grade-8/04-searching-and-sorting/outputs/binary_search.txt}{Never
more than 7 guesses between 1 and 100.}\label{out:g8:searching-and-sorting:binary}

\code{//} divides and drops the remainder, so \code{middle} is always a whole
position. The trace below follows the search for 77 in a \omterm{def:g8:searching-and-sorting:sorted}{sorted} row of 16
numbers:

\omcode{code/grade-8/04-searching-and-sorting/listings/binary_search_trace.py}{1}{15}{The
search, with \code{low}, \code{middle} and \code{high} printed at each
step.}\label{lst:g8:searching-and-sorting:trace}

\Needspace*{9\baselineskip}
\omoutput{code/grade-8/04-searching-and-sorting/outputs/binary_search_trace.txt}{Four
steps for 16 numbers.}\label{out:g8:searching-and-sorting:trace}

\begin{omfigure}
\begin{tikzpicture}[x=0.72cm, y=0.8cm, font=\scriptsize]
  \foreach \v [count=\i from 0] in {2, 5, 8, 12, 16, 23, 38, 41, 47, 56, 61, 72, 77, 85, 90, 98} {
    \draw[black!60, fill=ocBlue!8] (\i,0) rectangle ++(0.72,0.7);
    \node at (\i+0.36,0.35) {\v};
  }
  \draw[ocMint!80!black, very thick, fill=ocMint!25] (12,0) rectangle ++(0.72,0.7);
  \node at (12.36,0.35) {77};
  \foreach \lo/\hi/\mi/\y/\s in {0/15/7/-0.5/1, 8/15/11/-1.1/2, 12/15/13/-1.7/3, 12/12/12/-2.3/4} {
    \draw[ocInk, thick] (\lo+0.05,\y+0.15) -- (\lo+0.05,\y) -- (\hi+0.67,\y) -- (\hi+0.67,\y+0.15);
    \fill[ocRed] (\mi+0.36,\y) circle (2.2pt);
    \node[anchor=east] at (-0.1,\y) {step \s};
  }
\end{tikzpicture}

{\small \omterm{def:g8:searching-and-sorting:binary-search}{Binary search} for 77 among 16 \omterm{def:g8:searching-and-sorting:sorted}{sorted} numbers: each bracket is the part
still searched, from \code{low} to \code{high}, and the red dot the middle
\omterm{def:g7:lists:list}{item} compared. Four steps.}
\end{omfigure}

How many guesses can \omterm{def:g8:searching-and-sorting:binary-search}{binary search} need? Each guess removes the middle \omterm{def:g7:lists:list}{item}
and at least half of the others. One guess is enough for 1 \omterm{def:g7:lists:list}{item}, two for 3,
three for 7: $k$ guesses are enough for $2^k - 1$ \omterm{def:g7:lists:list}{items}. With 100 \omterm{def:g7:lists:list}{items}, 6
guesses cover only $2^6 - 1 = 63$, but 7 cover $2^7 - 1 = 127$: the doublings
of \cref{ch:g8:binary-and-hexadecimal}, read backwards.

The halving only makes sense if the \omterm{def:g7:lists:list}{list} is \omterm{def:g8:searching-and-sorting:sorted}{sorted}. On an unsorted \omterm{def:g7:lists:list}{list}, the
\omterm{def:g4:programs-and-operating-system:program}{program} still runs, and still answers, but the answer can be wrong:

\Needspace*{9\baselineskip}
\omcode{code/grade-8/04-searching-and-sorting/listings/unsorted_binary.py}{1}{6}{\omterm{def:g8:searching-and-sorting:binary-search}{Binary
search} on a \omterm{def:g7:lists:list}{list} that is not \omterm{def:g8:searching-and-sorting:sorted}{sorted}.}\label{lst:g8:searching-and-sorting:unsorted}

\Needspace*{6\baselineskip}
\omoutput{code/grade-8/04-searching-and-sorting/outputs/unsorted_binary.txt}{15
is at position 2, but \omterm{def:g8:searching-and-sorting:binary-search}{binary search} says it is not
there.}\label{out:g8:searching-and-sorting:unsorted}

The first middle \omterm{def:g7:lists:list}{item} is 11, smaller than 15, so the search goes right, where
15 is not. No \omterm{def:g6:first-python-program:error}{error message} warns about it: making sure the \omterm{def:g7:lists:list}{list} is \omterm{def:g8:searching-and-sorting:sorted}{sorted} is
the programmer's job.

\section{Selection sort}

\begin{definition}[Selection sort]\label{def:g8:searching-and-sorting:selection-sort}
\emph{Selection sort}\index{selection sort} finds the smallest \omterm{def:g7:lists:list}{item} of a \omterm{def:g7:lists:list}{list}
and swaps it to the front, then does the same on the rest of the \omterm{def:g7:lists:list}{list}, one
position further each time.
\end{definition}

\omcode{code/grade-8/04-searching-and-sorting/listings/selection_sort.py}{1}{16}{\omterm{def:g8:searching-and-sorting:selection-sort}{Selection
sort}, counting its \omterm{def:g6:if-else:comparison}{comparisons}.}\label{lst:g8:searching-and-sorting:selection}

\Needspace*{5\baselineskip}
\omoutput{code/grade-8/04-searching-and-sorting/outputs/selection_sort.txt}{6
cards, 15 \omterm{def:g6:if-else:comparison}{comparisons}.}\label{out:g8:searching-and-sorting:selection}

\code{a, b = b, a} swaps two \omterm{def:g6:variables-and-input:value}{values} in one line. The \omterm{def:g7:functions:function}{function} changes the \omterm{def:g7:lists:list}{list}
it receives, \emph{in place}, as \code{append} does
(\cref{ch:g7:lists}), and returns its number of \omterm{def:g6:if-else:comparison}{comparisons}.

\begin{omfigure}
\begin{tikzpicture}[x=0.75cm, y=0.75cm, font=\small]
  \foreach \row/\vals/\done in {0/{7,3,9,1,8,4}/0, 1/{1,3,9,7,8,4}/1, 2/{1,3,9,7,8,4}/2, 3/{1,3,4,7,8,9}/3, 4/{1,3,4,7,8,9}/4, 5/{1,3,4,7,8,9}/6} {
    \foreach \v [count=\i from 0] in \vals {
      \pgfmathsetmacro{\shade}{\i < \done ? 1 : 0}
      \ifdim\shade pt>0.5pt
        \draw[black!60, fill=ocMint!30] (\i,-\row) rectangle ++(0.7,0.7);
      \else
        \draw[black!60, fill=white] (\i,-\row) rectangle ++(0.7,0.7);
      \fi
      \node at (\i+0.35,-\row+0.35) {\v};
    }
  }
  \foreach \row/\txt in {0/start, 1/{1 swapped with 7}, 2/{3 already in place}, 3/{4 swapped with 9}, 4/{7 already in place}, 5/{8 in place: sorted}} {
    \node[anchor=west, font=\footnotesize] at (6.1,-\row+0.35) {\txt};
  }
\end{tikzpicture}

{\small \omterm{def:g8:searching-and-sorting:selection-sort}{Selection sort} on 6 cards: after each pass, one more card (green) is in
its final place.}
\end{omfigure}

The first pass compares 5 pairs, the next 4, then 3, 2 and 1: $5+4+3+2+1 = 15$
\omterm{def:g6:if-else:comparison}{comparisons}, whatever the order of the cards.

\section{Insertion sort}

\begin{definition}[Insertion sort]\label{def:g8:searching-and-sorting:insertion-sort}
\emph{Insertion sort}\index{insertion sort} takes the \omterm{def:g7:lists:list}{items} one after the
other and slides each one back, among those already \omterm{def:g8:searching-and-sorting:sorted}{sorted} before it, until
it reaches its place.
\end{definition}

It is how many card players sort a hand: pick up the cards one at a time, and
put each one where it belongs among those already held.

\omcode{code/grade-8/04-searching-and-sorting/listings/insertion_sort.py}{1}{19}{\omterm{def:g8:searching-and-sorting:insertion-sort}{Insertion
sort}, counting its \omterm{def:g6:if-else:comparison}{comparisons}.}\label{lst:g8:searching-and-sorting:insertion}

\Needspace*{5\baselineskip}
\omoutput{code/grade-8/04-searching-and-sorting/outputs/insertion_sort.txt}{The
same 6 cards, 11 \omterm{def:g6:if-else:comparison}{comparisons}.}\label{out:g8:searching-and-sorting:insertion}

\begin{omfigure}
\begin{tikzpicture}[x=0.75cm, y=0.75cm, font=\small]
  \foreach \row/\vals/\done in {0/{7,3,9,1,8,4}/1, 1/{3,7,9,1,8,4}/2, 2/{3,7,9,1,8,4}/3, 3/{1,3,7,9,8,4}/4, 4/{1,3,7,8,9,4}/5, 5/{1,3,4,7,8,9}/6} {
    \foreach \v [count=\i from 0] in \vals {
      \pgfmathsetmacro{\shade}{\i < \done ? 1 : 0}
      \ifdim\shade pt>0.5pt
        \draw[black!60, fill=ocYellow!35] (\i,-\row) rectangle ++(0.7,0.7);
      \else
        \draw[black!60, fill=white] (\i,-\row) rectangle ++(0.7,0.7);
      \fi
      \node at (\i+0.35,-\row+0.35) {\v};
    }
  }
  \foreach \row/\txt in {0/start: 7 alone is sorted, 1/{3 slides before 7}, 2/{9 stays}, 3/{1 slides to the front}, 4/{8 slides before 9}, 5/{4 slides before 7: sorted}} {
    \node[anchor=west, font=\footnotesize] at (6.1,-\row+0.35) {\txt};
  }
\end{tikzpicture}

{\small \omterm{def:g8:searching-and-sorting:insertion-sort}{Insertion sort} on the same 6 cards: the yellow part is \omterm{def:g8:searching-and-sorting:sorted}{sorted}, and
grows by one card each pass.}
\end{omfigure}

\section{Counting comparisons to compare two methods}

Both sorts give the same \omterm{def:g8:searching-and-sorting:sorted}{sorted list}. Which one works less? The \omterm{def:g6:while-loops:counter}{counters}
answer, on \omterm{def:g7:lists:list}{lists} of 10, 100 and 1,000 numbers, already \omterm{def:g8:searching-and-sorting:sorted}{sorted}, in reverse, or
mixed. The mixed order is fixed in the \omterm{def:g4:programs-and-operating-system:program}{program}, so the counts are the same
at every run.

\omcode{code/grade-8/04-searching-and-sorting/listings/count_comparisons.py}{1}{25}{The
two sorts on the same \omterm{def:g7:lists:list}{lists}.}\label{lst:g8:searching-and-sorting:count}

\omoutput{code/grade-8/04-searching-and-sorting/outputs/count_comparisons.txt}{\omterm{def:g6:if-else:comparison}{Comparisons}
of the two sorts.}\label{out:g8:searching-and-sorting:count}

\omterm{def:g8:searching-and-sorting:selection-sort}{Selection sort} always makes the same number of \omterm{def:g6:if-else:comparison}{comparisons} for $n$ \omterm{def:g7:lists:list}{items}:
$(n-1) + (n-2) + \dots + 1$, which is $n(n-1)/2$, so 45 for 10 \omterm{def:g7:lists:list}{items} and
499{,}500 for 1,000. \omterm{def:g8:searching-and-sorting:insertion-sort}{Insertion sort} depends on the order. On a \omterm{def:g8:searching-and-sorting:sorted}{sorted list},
each \omterm{def:g7:lists:list}{item} is compared once with the one before it and stays: $n - 1$
\omterm{def:g6:if-else:comparison}{comparisons}. On a reversed \omterm{def:g7:lists:list}{list}, each \omterm{def:g7:lists:list}{item} slides all the way to the front,
and it makes as many \omterm{def:g6:if-else:comparison}{comparisons} as \omterm{def:g8:searching-and-sorting:selection-sort}{selection sort}. On the mixed \omterm{def:g7:lists:list}{lists} it
lies in between, about half of the most.

\begin{omfigure}
\begin{tikzpicture}
  \begin{axis}[width=8.5cm, height=5.6cm, xmin=0, xmax=205, ymin=0,
      xlabel={number of items $n$}, ylabel={comparisons},
      label style={font=\footnotesize}, tick label style={font=\footnotesize},
      scaled y ticks=false, yticklabel style={/pgf/number format/1000 sep={,}},
      clip=false,
      axis x line*=bottom, axis y line*=left]
    \addplot[thick, ocRed] table[x=n, y=selection] {code/grade-8/04-searching-and-sorting/figures/out/comparisons.dat};
    \node[anchor=west, font=\scriptsize, align=left, ocRed] at (axis cs:203,19900) {selection sort, any order;\\ insertion sort, reversed};
    \addplot[thick, ocBlue, dashed] table[x=n, y=insmixed] {code/grade-8/04-searching-and-sorting/figures/out/comparisons.dat};
    \node[anchor=west, font=\scriptsize, ocBlue] at (axis cs:203,10147) {insertion sort, mixed};
    \addplot[thick, ocMint!70!black] table[x=n, y=inssorted] {code/grade-8/04-searching-and-sorting/figures/out/comparisons.dat};
    \node[anchor=west, font=\scriptsize, ocMint!70!black] at (axis cs:203,600) {insertion sort, sorted};
  \end{axis}
\end{tikzpicture}

{\small \omterm{def:g6:if-else:comparison}{Comparisons} against the number of \omterm{def:g7:lists:list}{items}, counted by the chapter's
\omterm{def:g7:functions:function}{functions} for $n$ from 10 to 200: the top curve bends upwards, the bottom one
is almost flat.}
\end{omfigure}

When the \omterm{def:g7:lists:list}{list} doubles, $n(n-1)/2$ is about four times larger: the top curve
bends upwards. A \omterm{def:g8:searching-and-sorting:linear-search}{linear search} on a \omterm{def:g7:lists:list}{list} twice as long only works twice as
much. Book 2 makes this way of comparing methods precise.

\begin{remark}[Sorting in practice]\label{rem:g8:searching-and-sorting:sorted}
Python has a sort built in: \code{sorted(items)} returns a new \omterm{def:g8:searching-and-sorting:sorted}{sorted list},
and \code{items.sort()} sorts a \omterm{def:g7:lists:list}{list} in place. They use a cleverer method than
the two of this chapter, and are the tools to use in real \omterm{def:g4:programs-and-operating-system:program}{programs}. Writing
the two sorts by hand shows what such a tool does, and how to count its
work.
\end{remark}

\section{Exercises}

\begin{exercise}[$\star$]\label{exo:g8:searching-and-sorting:1}
\omterm{def:g8:searching-and-sorting:binary-search}{Binary search} looks for 37 in the \omterm{def:g8:searching-and-sorting:sorted}{sorted list}
\code{[3, 8, 15, 21, 29, 37, 44, 50, 61]}. Write \code{low}, \code{middle},
\code{high} and the middle \omterm{def:g7:lists:list}{item} at each step.
\end{exercise}

\begin{exercise}[$\star$]\label{exo:g8:searching-and-sorting:2}
Write the \omterm{def:g7:lists:list}{list} \code{[6, 2, 8, 5, 1]} after the first two passes of
\omterm{def:g8:searching-and-sorting:selection-sort}{selection sort}.
\end{exercise}

\begin{exercise}[$\star$]\label{exo:g8:searching-and-sorting:3}
Write the same \omterm{def:g7:lists:list}{list} \code{[6, 2, 8, 5, 1]} after the first two passes of
\omterm{def:g8:searching-and-sorting:insertion-sort}{insertion sort}.
\end{exercise}

\begin{exercise}[$\star$]\label{exo:g8:searching-and-sorting:4}
Write a \omterm{def:g7:functions:function}{function} \code{all_positions(items, target)} that returns the \omterm{def:g7:lists:list}{list} of
all the positions where \code{target} appears, and an empty \omterm{def:g7:lists:list}{list} if it does
not.
\end{exercise}

\begin{exercise}[$\star\star$]\label{exo:g8:searching-and-sorting:5}
This version of \omterm{def:g8:searching-and-sorting:binary-search}{binary search} never ends when it looks for 2 in
\code{[4, 9]}. Follow \code{low} and \code{high} by hand to see why, and fix
it.

\omcode{code/grade-8/04-searching-and-sorting/exercises/binary_bug.py}{1}{15}{A
\omterm{def:g8:searching-and-sorting:binary-search}{binary search} that can \omterm{def:g4:algorithms-repeat-and-decide:loop}{loop} forever.}\label{lst:g8:searching-and-sorting:binary-bug}
\end{exercise}

\begin{exercise}[$\star\star$]\label{exo:g8:searching-and-sorting:6}
This \omterm{def:g8:searching-and-sorting:insertion-sort}{insertion sort} prints \code{[5, 1, 3]}. Which \omterm{def:g7:lists:list}{item} does it forget, and
how do you fix it?

\omcode{code/grade-8/04-searching-and-sorting/exercises/insertion_bug.py}{1}{11}{An
\omterm{def:g8:searching-and-sorting:insertion-sort}{insertion sort} that misses an \omterm{def:g7:lists:list}{item}.}\label{lst:g8:searching-and-sorting:insertion-bug}
\end{exercise}

\begin{exercise}[$\star\star$]\label{exo:g8:searching-and-sorting:7}
Write a \omterm{def:g4:programs-and-operating-system:program}{program} that uses \code{binary_search} from
\cref{lst:g8:searching-and-sorting:binary} to find the largest number of
guesses it needs for a target between 1 and 1,000, trying every target.
Check the answer with powers of two.
\end{exercise}

\begin{exercise}[$\star\star$]\label{exo:g8:searching-and-sorting:8}
Write your own \omterm{def:g8:searching-and-sorting:selection-sort}{selection sort} that sorts a \omterm{def:g7:lists:list}{list} of words from the shortest to
the longest.
\end{exercise}

\begin{exercise}[$\star\star$]\label{exo:g8:searching-and-sorting:9}
Count by hand the \omterm{def:g6:if-else:comparison}{comparisons} of \omterm{def:g8:searching-and-sorting:selection-sort}{selection sort} on 5 \omterm{def:g7:lists:list}{items}, pass by pass.
Does it depend on the order of the \omterm{def:g7:lists:list}{items}?
\end{exercise}

\begin{exercise}[$\star\star\star$]\label{exo:g8:searching-and-sorting:10}
Find a \omterm{def:g7:lists:list}{list} of 3 numbers that contains 5 but on which \omterm{def:g8:searching-and-sorting:binary-search}{binary search} does not
find 5. Explain why \omterm{def:g8:searching-and-sorting:binary-search}{binary search} needs a \omterm{def:g8:searching-and-sorting:sorted}{sorted list}.
\end{exercise}

\begin{exercise}[$\star\star\star$]\label{exo:g8:searching-and-sorting:11}
Show that \omterm{def:g8:searching-and-sorting:insertion-sort}{insertion sort} makes exactly $n - 1$ \omterm{def:g6:if-else:comparison}{comparisons} on a \omterm{def:g8:searching-and-sorting:sorted}{sorted list} of
$n$ \omterm{def:g7:lists:list}{items}, and exactly $n(n-1)/2$ on a reversed one. Count the \omterm{def:g6:if-else:comparison}{comparisons} of
each pass, then add them up.
\end{exercise}

\begin{exercise}[$\star\star\star$]\label{exo:g8:searching-and-sorting:12}
Write a \omterm{def:g7:functions:function}{function} \code{sort_books(books)} that sorts the records of the town's
book table (\cref{ch:g8:files-and-data-tables}) by year, and by title for
books of the same year, with your own \omterm{def:g8:searching-and-sorting:insertion-sort}{insertion sort}.
\end{exercise}

\section{Problem: One Word among a Million}

\begin{problem}[{Weekend problem --- finding one word among a million}]\label{pb:g8:searching-and-sorting:1}
A \omterm{def:g8:dictionaries:dictionary}{dictionary} of a million words is kept as a \omterm{def:g8:searching-and-sorting:sorted}{sorted list}. How many
\omterm{def:g6:if-else:comparison}{comparisons} does it take to find a word, and is it worth sorting first?

\medskip\noindent\textbf{Part I --- By hand.}
\begin{enumerate}
  \item Between 1 and 16, guess the secret number 11 by halving: write each
    guess and the part left.
  \item Halve 1,000 again and again, keeping the larger half, until one \omterm{def:g7:lists:list}{item}
    is left. How many halvings is that?
  \item Why does halving help only when the \omterm{def:g7:lists:list}{list} is \omterm{def:g8:searching-and-sorting:sorted}{sorted}?
\end{enumerate}

\medskip\noindent\textbf{Part II --- The worst case, by program.}

\omcode{code/grade-8/04-searching-and-sorting/listings/worst_guesses.py}{1}{16}{The
largest number of guesses for \omterm{def:g7:lists:list}{lists} of 10 to a million
\omterm{def:g7:lists:list}{items}.}\label{lst:g8:searching-and-sorting:worst}

\begin{enumerate}[resume]
  \item Why does \code{worst_guesses} always set \code{low = middle + 1}?
  \item Run the \omterm{def:g4:programs-and-operating-system:program}{program}. How do the counts grow when the size is multiplied
    by 10?
  \item Explain the 20 for a million with a power of two.
  \item How many \omterm{def:g6:if-else:comparison}{comparisons} would a \omterm{def:g8:searching-and-sorting:linear-search}{linear search} need in the worst case for
    a million words?
\end{enumerate}

\medskip\noindent\textbf{Part III --- Sorting first.}

\omcode{code/grade-8/04-searching-and-sorting/listings/sort_or_scan.py}{1}{24}{1,000
searches without and with sorting first.}\label{lst:g8:searching-and-sorting:sort-or-scan}

\begin{enumerate}[resume]
  \item What does the \omterm{def:g4:programs-and-operating-system:program}{program} compare?
  \item Run it. What do the 1,000 \omterm{def:g8:searching-and-sorting:linear-search}{linear searches} cost, and what do sorting
    plus 1,000 \omterm{def:g8:searching-and-sorting:binary-search}{binary searches} cost?
  \item With \omterm{def:g8:searching-and-sorting:selection-sort}{selection sort}, is sorting first worth it for 1,000 searches?
    With \omterm{def:g8:searching-and-sorting:insertion-sort}{insertion sort}?
  \item What if the \omterm{def:g8:dictionaries:dictionary}{dictionary} is searched 100,000 times, and \omterm{def:g8:searching-and-sorting:sorted}{sorted} only once?
  \item State the result.
\end{enumerate}
\end{problem}
